\documentclass[10pt]{article}

\usepackage[utf8]{inputenc}

\usepackage[T1]{fontenc}

\usepackage{amsmath}

\usepackage{amsfonts}

\usepackage{amssymb}

\usepackage[version=4]{mhchem}

\usepackage{stmaryrd}

\usepackage{bbold}

\usepackage{graphicx}

\usepackage[export]{adjustbox}

\graphicspath{ {./images/} }



\begin{document}

Векторные дифференциальные операции:



\[

\operatorname{grad}_{\rho} f=\frac{\partial f}{\partial \rho}, \operatorname{grad}_{\varphi} f=\frac{1}{\rho} \frac{\partial f}{\partial \varphi} ;

\][



\textbackslash operatorname\{div\} \textbackslash boldsymbol\{a\}=\textbackslash frac\{1\}\{\textbackslash rho\} a\_\{\textbackslash rho\}+\textbackslash frac\{\textbackslash partial a\_\{\textbackslash rho\}\}\{\textbackslash partial \textbackslash rho\}+\textbackslash frac\{1\}\{\textbackslash rho\} \textbackslash frac\{\textbackslash partial a\_\{\textbackslash Phi\}\}\{\textbackslash partial \textbackslash varphi\} ;

]



$\Delta f=\frac{1}{\rho} \frac{\partial}{\partial \rho}\left(\rho \frac{\partial f}{\partial \rho}\right)+\frac{1}{\rho^{2}} \frac{\partial^{2} f}{\partial \varphi^{2}}=\frac{\partial^{2} f}{\partial \rho^{2}}+\frac{1}{\rho} \frac{\partial f}{\partial \rho}+\frac{1}{\rho^{2}} \frac{\partial^{2} f}{\partial \varphi^{2}}$.



О б о б ш е н ны м и П. к. наз. числа $r$ и $\psi$, связанные с декартовыми шрямоугольными координатами $x$ и $y$ формулами:



$x=a r \cos \psi, y=b r \sin \psi$,



где $0 \leqslant r<\infty, 0 \leqslant \psi<2 \pi, a>0, b>0, a \neq b$. Координатные линии: әллипсы ( $r=$ const) и лучи ( $\psi=$ const).



Д. Д. Соколов.



ПОМЕХОУСТоИчИВость - понятие, характеризующее способность системы передачи информации противостоять искажающему действию помех. Предельно достижимую П. для оптимального метода передачи часто наз. п о т е н ц и а л ь н о й П. В теории шередачи информации П. конкретной системы передачи информации характеризуют сообцений точностью воспроизведения и, в частности, ошибочного декодирования вероятностью переданного сообщения.



Лит.: [1] Х а р к е в и ч А. А., Борьба с помехами, 2 изд., M., 1965; [2] В о з е н к р а Ф т मे ж., Д же к о б c И., Теоретические основы техники связи, пер. с англ., М., 1969.



ПОНТРЯГИНА ДВОИСТВЕННОСТЬ - 1) П. д.двойственность между абелевыми топологич. группами и их характеров группами. ' Т е о р е м а д в о й с тв е н в о с т и утверждает, что если $G$ - локально компактная абелева группа и $X(G)$ - ее группа характеров, то естественный гомоморфизм $G \rightarrow X(X(G))$, переводящий $a \in G$ в характер $\omega_{a}: X(G) \rightarrow T$, заданный формулой



\[

\omega_{a}(\alpha)=\alpha(a), \alpha \in X(G),

\]



есть изоморфизм топологич. групп. Из этой теоремы выводятся следующие утверждения.



I. Если $H-$ замкнутая подгруппа в $G$ и



\[

H^{*}=\{\alpha \in X(G) \mid \alpha(H)=0\}

\]



\begin{itemize}

  \item ее аннулятор в $X(G)$, то $H$ совпадает с аннулятором $\left\{a \in G \mid \alpha(a)=0 \quad \forall \alpha \in H^{*}\right\}$

\end{itemize}



подгруппы $H^{*}$; при этом группа $X(H)$ естественно изоморфна $X(G) / H^{*}$, а $X(G / H)$ - группе $H^{*}$.



II. Если $\varphi: G \rightarrow H$ - непрерывный гомоморфизм локально компактных абелевых групп, то после отождествления группы $G$ с $X(X(G))$ и группы $H$ с $X(X(H))$ при помощи изоморфизмов $a \rightarrow \omega_{a}$ гомоморфизм $\left(\varphi^{*}\right)^{*}$ отождествляется с $\varphi$.



III. $B e c$ группы $X(G)$ (как топологич. пространства) совпадает с весом группы $G$.



П. д. сопоставляет компактным группам $G$ дискретные группы $X(G)$ и наоборот. При этом компактная группа $G$ метризуема тогда и только тогда, когда $X(G)$ счетна, и связна тогда и только тогда, когда группа $X(G)$ без кручения. Конечномерность компактной группы $G$ равносильна тому, что $X(G)$ имеет конечный ранг (см. Ранг абелевой группы). Конечномерная компактная группа $G$ локально связна тогда и только тогда, когда $X(G)$ конечно порождена. Если $G$ конечна, то П. д. совпадает с двойственностью конечных абелевых групп, рассматриваемой над полем $\mathbb{C}$ комплексных чисел.



Топологич. группы, для к-рых верна теорема двойственности, наз. р е ф л е с и в ны ы и. Они не исчерпываются локально компактными группами, т. к. любое банахово пространство, рассматриваемое как топо- [ логич. группа, рефллексивно [8]. 0 характеризации рефлексивных групп см. [9].



Аналог П. Д. известен и для некоммутативных групп (те о рем а д во й с т в е нно с ти Т а на к а К р е й н a) (см. [4], [6], [7]). Пусть $G$ - компактная топологич. группа, $R$ - алгебра комплекснозначных представляющих функций на $G, S(R)$ - множество всех ненулевых гомоморфизмов алгебр $\omega: R \rightarrow \mathbb{C}$, удовлетворяюгих условию $\omega(\overline{f)}=\overline{\omega(f)}, f \in R$. В $S(R)$ определяется умножение $(\alpha, \beta) \rightarrow \alpha \beta$, удовлетворяющее следующему условию: если $\rho: G \rightarrow U(m)$ - неприводимое непрерывное унитарное представление группы $G$ и $\rho_{i j}(g)$ - его матричные элементы, то



\[

\left\|(\alpha \beta)\left(\rho_{i j}\right)\right\|=\left\|\alpha\left(\rho_{i j}\right)\right\| \cdot\left\|\beta\left(\rho_{i j}\right)\right\| \text {. }

\]



Это умножение и естественная топология превращают $S(R)$ в топологич. группу. Каждому $g \in G$ отвечает гомоморфизм $\alpha_{g} \in S(R)$, задаваемый формулой



\[

\alpha_{g}(f)=f(g), f \in R .

\]



Тогда соответствие $g \rightarrow \alpha_{g}$ есть изоморфизм топологич. группы $G$ на $S(R)$. Дано также алгебраич. описание категории алгебр $R$, к-рая оказывается, таким образом, двойственной категории компактных топологич. групп. Эта теория допускает обобцение на случай однородных пространств компактных топологич. групп (см. [4]).



Jum.: [1] П о н т р я г и н J. С., "Ann. Math.", 1934, v. 35, № 2, p. 361-88 (рус. пер.- «Успехи матем. наук», 1936, в. 22, 1973; [3] K a m p e n E. v a n, "Ann. Math.», 1935, v. 36, p 448-63; [4] K р е й н м. Г., "У кр. матем. ж.", 1949, т. 1 , № 4 , c. $64-98 ; 1950$, т. 2 , № 1, с. $10-59 ;$ 5] М о p p и с С., Двойственность Понтрягина и строение локально компактных абелевых групп, пер. с англ., М., $1980 ; 16]$ Н а й м а р к М. А.,



\includegraphics[max width=\textwidth, center]{2023_07_08_b903d1e18190bdeab223g-1}

T. 2, M., 1975; [8] S m i t h M. F., "Ann. Math.", 1952, v. 56, 1976, Bd 149, H. 2, S. $109-19$. A. Л. Онищик.



\begin{enumerate}

  \setcounter{enumi}{1}

  \item П. Д. в то п о л о г и и - изоморфизм между $p$-мерной группой когомологий Александрова - Чеха $H^{p}(A ; G)$ с коэффициентами в группе $G$ компактного множества $A$, лежацего в $n$-мерном компактвом ориентируемом многообразии $M^{n}$, и $(n-p-1)$-мерной группой гомологий $H_{n-p-1}^{c}(B ; G)$ дополнения $B=M^{n} \backslash A$ в предположении, что $H^{p}\left(M^{n} ; G\right)=H^{p+1}\left(M^{n} ; G\right)=0$ (гомологии и когомологии в размерности нуль - приведенные; символ $c$ означает компактные носители). Для случая, когда $A$ или $B$ - конечный полиэдр, этот изоморфизм был установлен Дж. Александером (J. Alexander). H. Стинрод (N. Steenrod) установил наличие такого изоморфизма для любого открытого подмножества $A \subset M^{n}$, a К. А. Ситников - для произвольного подмножества $A$.

\end{enumerate}



В приведенном виде закон двойственности Понтрягина был сформулирован П. С. Александровым. В первоначальной форме утверждалась двойственность в смысле теорип характеров между группами $H_{p}\left(A ; G^{*}\right)$ и $H_{n-p-1}^{c}(B ; G)$, где $G^{*}$ - бикомпактная группа характеров дискретной группы $G$. Эквивалентность обеих формулировок закона двойственности следует из того, что группа $H_{p}\left(A ; G^{*}\right)$ есть группа характеров группы $H^{P}(A ; G)$. Из предположения об ацикличности многообразия в размерностях $p$ и $p+1$ и из точной последовательности когомологий пары $\left(M^{n}, A\right)$ вытекает, что $H^{p}(A ; G)=H^{p+1}\left(M^{n}, A ; G\right)$, поэтому П. Д.- простое следствие двойственности Пуанкаре - Лефшеца (см. Пуанкаре двойственность).



Наиболее общая форма соотношений двойственности рассматриваемого типа состоит в следующем. Пусть $M^{n}$ - произвольное многообразие (возможно, обобщенное, не обязательно компактное и не обязательно ориентируемое), $\quad \mathcal{-}$ - локально постоянная система





\end{document}